\subsection{Sprint 2}

A Sprint 2 teve como foco avançar na experiência de uso dos dois perfis centrais da plataforma: ONGs e empresas. As atividades previstas estavam concentradas principalmente no frontend, com a entrega da visualização principal dos projetos da ONG e a construção da vitrine de busca e modalidades de investimento para empresas. A primeira frente tinha como objetivo permitir que a ONG acompanhasse rapidamente seus projetos, visualizando informações relevantes sem precisar navegar por fluxos excessivamente longos. A segunda frente buscava melhorar a descoberta de oportunidades por parte das empresas, aproximando-as de projetos alinhados aos seus interesses.

Todas as atividades previstas para a sprint foram concluídas. Do ponto de vista do produto, essa etapa foi importante porque começou a tornar a plataforma mais utilizável para os perfis finais. Até então, boa parte do esforço estava concentrada na estruturação da base do sistema, autenticação, modelagem e fluxos iniciais. Na Sprint 2, o projeto passou a ganhar mais corpo visual e funcional, especialmente na experiência de navegação das ONGs e das empresas.

Minha atuação esteve relacionada ao acompanhamento das demandas do squad, ao entendimento das entregas previstas e ao apoio na organização das atividades de frontend. Como eu já possuía mais experiência do que em AGES anteriores, passei a perceber com mais clareza os impactos de uma user story mal definida ou de uma divisão de responsabilidades pouco explícita. Esse foi justamente um dos principais problemas encontrados durante a sprint.

A maior dificuldade foi a falta de clareza nas user stories e nas funções de cada squad. No início, entendíamos que o nosso squad, chamado de User Experience, teria como responsabilidade principal concluir tarefas relacionadas ao frontend. Entretanto, ao longo da sprint, ficou claro que a expectativa não era apenas implementar telas, mas também pensar fluxo, experiência, coerência do produto, alinhamento com regras de negócio e comunicação com outros squads. Essa diferença de entendimento gerou ruídos, pois algumas tarefas pareciam pertencer a uma frente, mas dependiam diretamente de decisões ou entregas de outra.

Esse problema evidenciou uma fragilidade importante na comunicação entre AGES III e AGES IV. Como os integrantes da AGES III estavam em posição de maior responsabilidade na gestão dos squads, era esperado que tivéssemos uma postura mais ativa na interpretação das demandas, na divisão das tarefas e na comunicação com os demais membros. No entanto, percebi que ainda havia dificuldade em transformar decisões de alto nível em tarefas claras para execução.

A principal lição aprendida foi que a comunicação entre AGES III e AGES IV precisava melhorar, especialmente porque os AGES III estavam cada vez mais envolvidos na gestão dos squads. Liderar tecnicamente ou organizar um grupo não significa apenas distribuir tarefas, mas garantir que todos entendam o propósito da entrega, os critérios de aceite e as dependências envolvidas. Quando esse alinhamento não acontece, mesmo tarefas aparentemente simples podem gerar retrabalho ou insegurança na equipe.

Como próximo passo, entendi que seria necessário desenvolver o projeto de forma mais próxima dos AGES IV, buscando validar expectativas com antecedência e reduzir ambiguidades. Essa sprint reforçou que, em projetos colaborativos, a qualidade da entrega depende tanto da implementação quanto da clareza do processo que leva até ela.
